\chapter{Introduction}
\label{ch:introduction}

Write your introduction here. A good introduction motivates the problem, states the research questions you aim to answer, and closes with an outline of the thesis structure. Back up important claims with citations.\footnote{A Bachelor's thesis is typically around 50 pages; a Master's thesis around 80 pages. The exact length depends on your programme and the scope of your research.}

%%%%%%%%%%%%%%%%%%%%%%%%%%%%%%%%%%%%%%%%%%%%%%%%%%%%%%%%%%%%%%%%%%%%%%%%
\section{Motivation}
\label{sec:motivation}

Describe the broader context and explain why the problem you are tackling is relevant. What gap in existing work does your thesis address? Why should the reader care?\footnote{Use citations to ground your motivation in the literature. Unsupported claims weaken the introduction.}

%%%%%%%%%%%%%%%%%%%%%%%%%%%%%%%%%%%%%%%%%%%%%%%%%%%%%%%%%%%%%%%%%%%%%%%%
\section{Research Questions}
\label{sec:research-questions}

State the specific research questions or objectives of your thesis. For example:

\begin{itemize}
  \item[\textit{RQ1}] Your first research question.
  \item[\textit{RQ2}] Your second research question.
  \item[\textit{RQ3}] Your third research question.
\end{itemize}

%%%%%%%%%%%%%%%%%%%%%%%%%%%%%%%%%%%%%%%%%%%%%%%%%%%%%%%%%%%%%%%%%%%%%%%%
\section{Outline}
\label{sec:outline}

Briefly describe what each chapter covers. For example:

Chapter~\ref{ch:relatedwork} surveys the existing literature and positions this work within it.
Chapter~\ref{ch:ownwork} presents the approach and its implementation.
Chapter~\ref{ch:evaluation} evaluates the approach against the research questions.
Chapter~\ref{ch:summary} summarises the contributions and outlines directions for future work.
